\section{CONCLUSIONS}

In this review, we presented a unified, PRISMA-grounded overview of optical integrated sensing and communication (O-ISAC) and positioned the field as a family of related architectures rather than as a single interchangeable design space. The central motivation throughout the manuscript has been that optical carriers open a materially different operating regime for joint communication and sensing, but that the resulting literature cannot be synthesized credibly unless modality, metric, and deployment context are kept explicit.

Specifically, we first established the physical-layer foundations that motivate O-ISAC for 6G-facing systems, including the bandwidth, resolution, and hardware arguments that make optical carriers attractive beyond conventional RF-limited designs. We then consolidated a PRISMA 2020-compliant evidence base of 220 studies and used that corpus to build a cross-modality taxonomy spanning fiber, free-space optics, VLC/LiFi, photo-THz-linked, and hybrid settings. This taxonomy was not introduced as a naming exercise alone; it provides the structural basis for comparing studies only when medium, integration style, detection regime, and task scope are aligned.

On top of that structural layer, the manuscript introduced a reporting and comparison contract for performance synthesis. A key conclusion of this review is that bandwidth-limited physical resolution, estimator-level accuracy, and information-theoretic bounds should not be collapsed into a single generic ``sensing performance'' label, and signal-quality quantities should not be mixed across optical and electrical measurement planes without explicit justification. Under this governed reading, Section V showed that admissible rate-resolution evidence remains sparse and unevenly distributed, Section VI showed that enabling stacks such as machine learning, optical RIS, optical phased arrays, and photonics-assisted high-frequency generation are promising but not yet sufficient to erase reporting and benchmarking fragmentation, and Section VII showed that deployment evidence is concentrated in a limited set of application settings rather than being broadly mature across all domains.

These observations motivate the roadmap synthesized in Section VIII. The next phase of O-ISAC research should not prioritize isolated record demonstrations alone; it should prioritize interoperable reporting practice, reproducible evaluation workflows, clearer channel and benchmark definitions, implementation-scalable hardware pathways, and domain-aware deployment validation. In other words, the field needs not only better devices and higher headline performance, but also stronger comparability discipline and more transparent evidence packaging.

Taken together, this review aims to serve as both a synthesis and a guardrail. It provides a traceable map of what has been demonstrated, where cross-domain transfer appears plausible, and why several widely cited comparisons remain methodologically unsafe when metric governance is ignored. We hope this manuscript helps researchers frame O-ISAC as a rigorous systems problem whose long-term impact will depend as much on methodological consistency and deployment realism as on photonic innovation itself.
